\section{Smooth arithmetic families from the original smooth fibre}

\begin{lemma}[Smoothness descends on the actual residue fibre]
\label{lem:kn-arithmetic-original-smooth-fibre}
\leanok
Suppose a square with upper scheme $X'$, lower scheme $T$, right
morphism $g:Y\to S$, and bottom morphism $T\to S$ is cartesian.
If the original upper structural morphism $X'\to T$ is smooth, then,
for every $t\in T$, the actual fibre $Y_{h(t)}$ is smooth over its
original residue field $\kappa(h(t))$.
\end{lemma}
\begin{proof}
Take the native residue fibres of the whole cartesian square.  Their
square is again cartesian, with bottom arrow
$\operatorname{Spec}\kappa(t)\to\operatorname{Spec}\kappa(h(t))$.
A homomorphism between fields is faithfully flat, and its spectrum is
quasicompact and surjective.  The original smooth morphism has smooth
residue fibre by base change.  Faithfully flat descent of smoothness
in this actual fibre square proves the assertion.  No generic-fibre
smoothness hypothesis is supplied separately.
\end{proof}

\begin{lemma}[The original standard chart produces its dimension-preserving model]
\label{lem:kn-arithmetic-original-dimension-chart-model}
\lean{EnvelopeProperty.Vanishing.ArithmeticOriginalSmooth.exists_original_standard_chart_model}
\leanok
Let the original morphism to the coefficient-system limit be smooth of
relative dimension $n$. Every original point has an actual affine open
neighbourhood with a finite-presentation coefficient model smooth of
that same relative dimension $n$.
\end{lemma}
\begin{proof}
Use the actual standard-smooth chart of relative dimension $n$ over the
whole original affine base, with its original coefficient map and
structural triangle. Native Noetherian descent for standard-smooth
algebras retains $n$ and produces a finite coefficient algebra and a
whole tensor equivalence. The spectrum of this equivalence is the
actual cartesian model. The original chart isomorphism transports
that model back to the unchanged original affine neighbourhood.
\end{proof}

\begin{theorem}[The same original fixed family retains its relative dimension]
\label{thm:kn-arithmetic-original-dimension-spread}
\lean{EnvelopeProperty.Vanishing.ArithmeticOriginalSmooth.exists_relativeDimension_fixed_model_baseChange}
\leanok
\uses{lem:kn-arithmetic-original-dimension-chart-model}
Suppose the original morphism is smooth of relative dimension $n$ and
has a fixed finite-presentation model in the actual finite-coefficient
approximation system. At some produced later coefficient stage, the
base change of this same fixed model is smooth of relative dimension
$n$.
\end{theorem}
\begin{proof}
The original scheme is quasicompact because it is the affine-system
limit of base changes of its finite-presentation model. Take a finite
subcover of the actual standard charts of the preceding lemma.
Native approximation descends each chart to an actual open of a later
base change of the fixed original model. The original chart model and
this descended open become isomorphic after a further coefficient
extension, so the open has the same relative-dimension-$n$ property.

Choose a common later stage for the finitely many chart stages.
Relative-dimension-$n$ smoothness persists on each inverse-image open
by actual base change. Their union pulls back to the whole original
scheme. The native finite-presentation limit theorem then gives a
further stage where the union is the entire actual model. Since
smoothness of relative dimension $n$ is local on the source, those
opens prove the asserted property for that same model at the produced
stage. No smooth-model or relative-dimension certificate is supplied
as a final vanishing hypothesis.
\end{proof}

\begin{theorem}[The original prescribed dimension on the actual arithmetic family]
\label{thm:kn-arithmetic-original-dimension-family}
\lean{EnvelopeProperty.Vanishing.ArithmeticOriginalSmooth.exists_relativeDimension_smooth_projective_arithmetic_family,EnvelopeProperty.Vanishing.ArithmeticOriginalSmooth.isPullback_scalarTransition}
\leanok
\uses{thm:kn-arithmetic-original-dimension-spread}
The actual smooth projective arithmetic family of the original closed
embedding can be produced with the original prescribed relative
dimension $n$. Its arithmetic base is smooth and finite type over
$\mathbb Z$, its original field fibre is the whole original $X$ with
the original structural morphism, and all original projective twists
are recovered as whole sheaves.
\end{theorem}
\begin{proof}
Apply fixed-model dimension spreading to the actual original model
constructed earlier, preserving its original top and bottom coefficient
isomorphisms. The produced stage is a finite-type integer subalgebra
of the original field. Its canonical pullback family has the
prescribed dimension, and the original whole cartesian square is
unchanged. Localize the actual arithmetic stage on a produced
nonzero smooth principal open. The canonical coefficient transition
is itself an actual pullback square, so dimension-$n$ smoothness
persists. The already proved original embedding and projective-twist
comparisons apply to this same enlarged coefficient subalgebra.
\end{proof}

\begin{theorem}[Actual dimension-preserving Witt and residue families]
\label{thm:kn-arithmetic-original-dimension-witt}
\lean{EnvelopeProperty.Vanishing.ArithmeticOriginalSmooth.exists_relativeDimension_witt_stage,EnvelopeProperty.Vanishing.ArithmeticOriginalSmooth.exists_relativeDimension_witt_stage_over,EnvelopeProperty.Vanishing.ArithmeticOriginalSmooth.residueSpecialFibreIso_hom_structure}
\leanok
\uses{thm:kn-arithmetic-original-dimension-family}
For any prescribed characteristic bound at least $n$, the same
original family produces an actual closed residue field of prime
characteristic $p>n$ and an actual smooth length-two Witt family of
relative dimension $n$. The actual residue family is canonically
isomorphic, over that same residue field, to the literal special
fibre used by the temporary DI input.
\end{theorem}
\begin{proof}
Apply the actual arithmetic closed-point/Witt coefficient lift to
the dimension-preserving smooth arithmetic base. Its actual further
localization preserves the family and its relative dimension by the
canonical coefficient square. Base change by the produced whole
Witt coefficient homomorphism gives the actual Witt family; base
change by its actual residue quotient gives the actual residue
family. Smoothness of the same prescribed dimension survives both
base changes, and the produced prime exceeds the chosen bound.

The whole Witt reduction equality gives the existing actual residue
family isomorphism to the pullback in the opposite factor order.
Compose it with the canonical pullback symmetry. Its first projection
is exactly the original residue structural morphism, so it identifies
the actual family with the literal DI special fibre over the
unchanged residue field. No supplied fibre-comparison certificate
is used.

If original-line descent has already enlarged the coefficient stage,
start with that actual finite-type stage and its proved
relative-dimension property. The same localization and closed-point
construction retains the entire given stage in the produced Witt
base. Thus all original-line coefficients are fixed before choosing
the arithmetic closed point; smoothness of that enlarged integer
base is produced by localization.
\end{proof}

Earlier constructions supply the actual projective arithmetic and Witt-family
constructions reused here. Their model twists recover
$j^*\mathcal O(d)$. For an arbitrary ample $L$, a very ample power
$L^r$ does not by itself descend $L$; a separate construction
provides the required actual line descent.
